\SourceLeaf{320}%
zum Grunde lag, welcher in der gegenwärtigen Untersuchung genügt ist. Man hat dabei nur zu erwägen, dass die dortigen \(\psi^0\) und \(\psi'\) nichts anderes sind, als hier \(T+\frac12t-(U+\frac12u)\) und \(T-\frac12t-(U-\frac12u)\); ferner das dortige \(h\) dasselbe was hier \(s\); endlich dass das dortige \(\chi\) von der hier mit \(U\) bezeichneten Grösse im Allgemeinen nur um eine Grösse zweiter Ordnung verschieden sein kann, jedenfalls aber der Unterschied wenigstens von der ersten Ordnung ist. Es ergibt sich so, auf die dritte Ordnung einschl. genau
\[
\begin{aligned}
T+\tfrac12t&=U+\tfrac12u-\frac{ee\cos P\mathord{.}\sin P\mathord{.}\sin U\mathord{.}\cos U^2}{12\cos\varphi\cos\theta}\mathord{.}s^3\\
T-\tfrac12t&=U-\tfrac12u+\frac{ee\cos P\mathord{.}\sin P\mathord{.}\sin U\mathord{.}\cos U^2}{12\cos\varphi\cos\theta}\mathord{.}s^3
\end{aligned}
\]
und folglich, eben so genau,
\[
\begin{aligned}
T&=U\quad .\ .\ .\ .\ .\ .\ .\ .\ .\ .\ .\ .\ .\ .\ .\ .\ .\ .\ .\ .\quad(6)\\
t&=u-\frac{ee\cos P\mathord{.}\sin P\mathord{.}\sin U\mathord{.}\cos U^2}{6\cos\varphi\cos\theta}\mathord{.}s^3\quad .\ .\ .\ .\ .\ .\ .\ .\ .\quad(7)
\end{aligned}
\]

\indent Die Vergleichung der Länge der geodaetischen Linie auf dem Ellipsoid mit dem Grösstekreisbogen auf der Kugel ist zwar in Art. 15 für den in Rede stehenden Fall nicht besonders entwickelt: es ist jedoch sehr leicht, diess zu ergänzen. Es ist nemlich in den dortigen Bezeichnungen die Länge des geodaetischen Bogens
\[
=A\int\frac{\cos y}{m\cos\psi}\mathord{.}\mathrm{d}x
\]
welche Integration von \(x=-\frac12(h-\delta)\) bis \(x=+\frac12(h+\delta)\) auszudehnen ist. Da \(y\) und \(\psi\) nur Grössen von der dritten Ordnung sind, so sieht man leicht, dass die Weglassung des Factors \(\frac{\cos y}{\cos\psi}\) in dem Werthe des Integrals nur einen Fehler der siebenten Ordnung hervorbringen kann. Jene Länge ist also, bis auf die fünfte Ordnung einschl. genau,
\[
\begin{aligned}
=A\int\frac{\mathrm{d}x}{m}&=A\int\mathrm{d}x(1-\mu x^3-\mu'x^4)\\
&=A(x-\tfrac14\mu x^4-\tfrac15\mu'x^5)+\text{Const.}\\
&=A\{h-\tfrac18(h^3\delta+h\delta^3)\mu-\tfrac1{80}(h^5+10h^3\delta\delta+5h\delta^4)\mu'\}
\end{aligned}
\]
Die Coëfficienten \(\mu,\mu'\) lassen sich angeben, wenn man in der Reihe
\[
m=1-\frac{2ee\cos P\mathord{.}\sin P}{3\cos\varphi\cos\theta}\mathord{.}q^3-\frac{ee\cos P^2}{6\cos\varphi^2\cos\theta^2}(1-7ee\sin P^2)q^4\quad .\ .\ .\ .
\]
(welche von selbst aus der Art 9 gegebenen folgt) für \(q\) die Substitution macht
